\chapter{Formulação do problema}

Nesta pesquisa de Mestrado será simulada computacionalmente a dinâmica orbital e de atitude de uma espaçonave equipada com manipulador robótico.
Ela está inserida em uma missão de \emph{rendezvous and docking/berthing} (RVD/B) em órbita baixa terrestre (LEO, \emph{Low Earth Orbit}).
A formulação contempla duas camadas principais:
(i) a dinâmica orbital relativa entre o \emph{chaser} e o \emph{target}, e é descrita pelas Equações de Hill/Clohessy-Wiltshire em coordenadas LVLH; e
(ii) a dinâmica de atitude da base flutuante e do manipulador robótico acoplado, e são modelados segundo a abordagem de corpos rígidos interconectados.

No subsistema orbital, considera-se que a aproximação é realizada pela direção radial (\emph{R-bar}), e que a transferência inicial entre órbitas circulares é obtida por meio de uma manobra de Hohmann, sem perturbações ambientais.
Para o subsistema de atitude, são analisadas as equações de movimento da espaçonave com base flutuante, sujeita à ação dos torques de controle, e as equações de movimento das juntas do manipulador, composto por quatro articulações revolutas e uma prismática (configuração $4R1P$), responsáveis pelo alinhamento e extensão da ferramenta de captura.

A estratégia de controle empregada é realizada através do uso de controladores do tipo proporcional$-$derivativo (PD).
Estes são aplicados tanto à atitude da base quanto às juntas do manipulador, assegurando o alinhamento do \emph{chaser} em relação ao \emph{target}.
A integração entre a dinâmica orbital, o controle de atitude da espaçonave e o controle do manipulador robótico é realizada em ambiente de microgravidade.
Para garantir a fidelidade numérica das simulações, adota-se a consideração integral dos produtos de inércia, mesmo quando seus valores são da ordem de $10^{-6}$ ou inferiores.
Dessa forma, a simulação proposta abrange, de maneira unificada, os aspectos orbitais (transferência e aproximação pelo \emph{R-bar}), os aspectos de atitude (base livre em microgravidade), e os aspectos robóticos (cinemática e dinâmica de um manipulador $\mathbf{4R1P}$ controlado por controle PD), fornecendo uma plataforma para análise de desempenho e robustez de missões de RVD/B.

\section{Hipóteses de modelagem}
A fase de aproximação considerada neste trabalho corresponde à etapa transitória do \emph{rendezvous and docking/berthing}, anterior à captura pela ferramenta (\emph{tool capture}), momento em que o \emph{chaser} permanece em espera em relação ao \emph{target}.
Após o acoplamento da ferramenta à porta de docking, as propriedades dinâmicas dos sistemas \emph{chaser} e \emph{target} se modificam, passando a constituir um único corpo rígido, a dinâmica desse novo sistema acoplado não é abordada neste estudo.

As espaçonaves são modeladas em regime de microgravidade, ou seja, ocorre o equilíbrio entre a força gravitacional e a força centrífuga, e não à ausência total de gravidade.
No subsistema do manipulador robótico, as forças gravitacionais são desprezadas; já no subsistema orbital relativo, a translação e a rotação relativas são descritas pelas Equações de Hill \cite{hill1878}. Dessa forma, este documento trata apenas da atitude da base e do movimento do manipulador robótico.
Portanto, as Equações de Hill são utilizadas para representar o movimento orbital relativo entre as espaçonaves, assumindo órbita circular.

De forma complementar, o movimento orbital é assumido como Kepleriano, sem perturbações ambientais.
Por conseguinte, devido ao curto intervalo de tempo de simulação, essas perturbações não influenciam de forma significativa na modelagem.
Dessa forma, todos os corpos (juntas, elos, espaçonaves, base e porta de docking) são rígidos durante a fase não operacional, enquanto que na fase operacional, o manipulador robótico é tratado como corpo não-rígido.  

Em síntese, os efeitos como a flexibilidade estrutural, deformações locais, vibrações, folgas e aquecimento por atrito não foram incluídos na modelagem.
A formulação cinemática e dinâmica segue Craig \cite{craig2005}, e no contexto de microgravidade e missões de RVD/B com base flutuante são adotadas as referências de \cite{fehse2003,wie2008space}.

Por fim, as massas do \emph{chaser}, do \emph{target} e dos elos do manipulador robótico são consideradas constantes, sem consumo de propelente ou variação de massa.

\subsection{Equações do Movimento de Atitude em Missões de RVD/B}
O controle de atitude é utilizado em missões de \emph{rendezvous and docking/berthing}, pois permite que a orientação do \emph{chaser} permaneça dentro de limites pré-estabelecidos ao longo da aproximação.
Essa etapa é essencial nas manobras realizadas no eixo radial (\emph{R-bar}), em que o alinhamento entre as espaçonaves é condição para o sucesso da missão.

As equações de movimento de atitude podem ser expressas em diferentes sistemas de coordenadas: o inercial, o orbital (\emph{Local-Vertical-Local-Horizontal, LVLH}) e o relativo.
Cada sistema de referência é empregado conforme a necessidade de modelagem:
a referência inercial descreve a dinâmica absoluta da espaçonave;
o LVLH é usado para demonstrar a aproximação orbital relativa;
e o sistema relativo é útil para análise do alinhamento entre \emph{chaser} e \emph{target}.

De modo adicional, adota-se a hipótese de considerar a totalidade dos produtos de inércia, para assegurar fidelidade nas simulações dinâmicas, mesmo quando seus valores são muito pequenos (da ordem de $10^{-6}$ ou inferiores).
Por fim, os requisitos da missão estabelecem que a atitude do \emph{chaser} deve permanecer alinhada ao \emph{target} ao longo do \emph{R-bar}, respeitando as restrições estabelecidas sobre a velocidade angular.

\subsection{Atuadores e sensores}
Os atuadores das juntas revolutas aplicam torques, enquanto o atuador da junta prismática exerce força.
Ambos são representados por $\tau_i$, e a modelagem é abstrata no nível dinâmico, sem referência direta ao hardware.

\subsection{Restrições operacionais}
Para evitar colisões, o movimento do manipulador robótico é restringido a um espaço de trabalho (\emph{workspace}) previamente definido, de modo a não interceptar os limites físicos do \emph{chaser} ou do \emph{target}.
Além disso, cada junta está sujeita a limites de posição, velocidade e aceleração.
Outrossim, as quatro primeiras juntas ($J_1 - J_4$) são do tipo revolutas e a quinta ($J_5$) é do tipo prismática, resultando na configuração $4R1P$, com curso limitado ao seu comprimento útil.

% Além disso, cada junta está sujeita a limites de posição \eqref{eq:limites_juntas_posicao}, velocidade \eqref{eq:limites_juntas_velocidade} e aceleração \eqref{eq:limites_juntas_aceleracao}:
% \begin{equation}
%     \label{eq:limites_juntas_posicao}
%     q_{i,\min} \le q_i \le q_{i,\max}
% \end{equation}

% \begin{equation}
%     \label{eq:limites_juntas_velocidade}
%     |\dot q_i| \le \dot q_{i,\max}
% \end{equation}

% \begin{equation}
%     \label{eq:limites_juntas_aceleracao}
%     |\ddot q_i| \le \ddot q_{i,\max}
% \end{equation}

\subsection{Incertezas e variações}
Foram consideradas incertezas paramétricas fixas associadas a cada elo do manipulador e ao conjunto, abrangendo as massas $m_i$, as posições dos centros de massa (CoM) e os momentos de inércia.
Tais incertezas foram modeladas como desvios percentuais em torno dos valores nominais, com limites máximos pré-definidos.
Cada simulação é realizada com um conjunto fixo de parâmetros perturbados dentro desses limites, representando erros de modelagem ou variações construtivas.  

As variáveis de estado adotadas no modelo dinâmico são os ângulos (ou deslocamento prismático) das juntas e suas velocidades generalizadas.
As acelerações são obtidas diretamente das equações de movimento e, portanto, não constituem estados independentes, mas variáveis derivadas.

\subsection{Condições iniciais}
As condições iniciais do sistema são definidas pelas posições no sistema de coordenadas local, pelas velocidades lineares e angulares iniciais da base, além dos ângulos iniciais de cada junta do manipulador e suas respectivas velocidades iniciais.

\subsection{Hipóteses de simulação e critérios de validação}

As simulações numéricas serão utilizadas como ferramenta para validar a transferência de Hohmann, as equações do movimento absoluto, do movimento relativo e a lei de controle proposta.
Tais simulações serão conduzidas em ambiente MATLAB/Simulink, permitindo avaliar a resposta dinâmica do sistema frente a condições iniciais estabelecidas e a diferentes níveis de complexidade na modelagem.

As hipóteses de configuração inicial incluem a definição das condições de atitude tanto do \emph{chaser} quanto do \emph{target}.
O \emph{chaser} possuirá desvios em sua posição e em suas velocidades iniciais, enquanto que o \emph{target} será tratado como referência, servindo como orientação fixa na referência local LVLH.
A integração temporal das equações diferenciais será realizada por métodos numéricos de passo variável na propagação dos estados e estabilidade no tratamento de equações não lineares.  

Serão avaliados diferentes cenários de simulação, abrangendo:
(i) a transferência orbital do \emph{chaser} para a órbita do \emph{target} por meio da transferência de Hohmann;
(ii) o controle absoluto de atitude na referência inercial;
(iii) o controle relativo de atitude no cenário entre \emph{chaser} e \emph{target}; e
(iv) a comparação de erros residuais entre as atitudes simuladas e as atitudes de referência.

Os critérios de avaliação adotados são:
(i) a estabilização, entendida como a convergência da atitude para a orientação desejada sem oscilações não físicas;
(ii) a robustez, medida pela manutenção de desempenho aceitável frente a variações inerciais do manipulador e a perturbações externas;
(iii) a saturação de atuadores, observando se os torques exigidos permanecem dentro dos limites físicos do sistema;
(iv) o tempo de resposta, avaliado pela duração necessária para que o erro angular relativo atenda aos critérios de desempenho; e
(v) o atendimento às condições de sincronização posição--atitude, em especial a exigência de erro angular inferior a $1^\circ$ e a permanência do \emph{chaser} dentro do corredor seguro de aproximação.
